% ch30-appendix-diagrams — ملحق الديجرامات الرسمية (embedded من official-diagrams.pdf)
\section{ملحق: الديجرامات الرسميّة للمشروع}

الديجرامات دي هي المخطّطات الرسميّة المعتمَدة للمشروع، مُدمَجة هنا بجودة \en{vector}
(تقدر تكبّرها في العارض من غير ما تتكسّر). كلّ ديجرام في صفحة عريضة \en{(landscape)}
لوحده. الترتيب: مخطّطات العلاقات \en{(ERD)}، ثمّ المُحسَّنة \en{(EERD)}، ثمّ الربط
العلائقيّ، ثمّ الأصناف \en{(Class)}، ثمّ المكوّنات والنشر \en{(Component)}. في فصل
«خريطة الكود للمناقشة» بنقول لك أي ديجرام تعرضه مع أي موضوع.

% مولّد صفحة ديجرام: #1 رقم الصفحة في official-diagrams.pdf ، #2 العنوان العربي
\newcommand{\diag}[2]{%
  \begin{landscape}%
  \begin{center}%
  {\small\bfseries\color{spPrimary} #2}\par\vspace{3pt}%
  \includegraphics[page=#1,width=\linewidth,height=0.86\textheight,keepaspectratio]{official-diagrams.pdf}%
  \end{center}%
  \end{landscape}}

\subsection{مخطّطات العلاقات (\en{ERD})}
\diag{1}{\en{ERD} --- كلّ الموديولات (نظرة شاملة)}
\diag{2}{\en{ERD} --- موديول 1: الهوية والوصول والملف الشخصيّ}
\diag{3}{\en{ERD} --- موديول 2: المنح والطلبات والمستندات}
\diag{4}{\en{ERD} --- موديول 3: حجز الاستشاريين والتقييمات}
\diag{5}{\en{ERD} --- موديول 4: المجتمع والدردشة}
\diag{6}{\en{ERD} --- موديول 5: مركز الموارد والإشعارات}
\diag{7}{\en{ERD} --- موديول 6: الذكاء الاصطناعيّ وقاعدة المعرفة والمنصّة}

\subsection{المخطّطات المُحسَّنة (\en{EERD})}
\diag{8}{\en{EERD} --- كلّ الموديولات (نظرة شاملة)}
\diag{9}{\en{EERD} --- تخصيص المستخدم (\en{USER Specialization}): طالب / استشاريّ / جهة مانحة / أدمن}

\subsection{الربط العلائقيّ (\en{Relational Mapping})}
\diag{16}{الربط العلائقيّ --- كلّ الموديولات (تحويل الكيانات إلى جداول)}

\subsection{مخطّطات الأصناف (\en{Class Diagrams})}
\diag{24}{مخطّط الأصناف --- موديول 1: الهوية والملف الشخصيّ}
\diag{25}{مخطّط الأصناف --- موديول 2: المنح والطلبات والمستندات}
\diag{26}{مخطّط الأصناف --- موديول 3: الحجز والاستشارات والتقييمات}
\diag{27}{مخطّط الأصناف --- موديول 4: المجتمع والدردشة}
\diag{28}{مخطّط الأصناف --- موديول 5: مركز الموارد والإشعارات}
\diag{29}{مخطّط الأصناف --- موديول 6: الذكاء الاصطناعيّ والمعرفة والمنصّة}

\subsection{مخطّطات المكوّنات والنشر (\en{Component \& Deployment})}
\diag{30}{مخطّط المكوّنات والنشر الرئيسيّ (\en{free-mode})}
\diag{31}{مخطّط المكوّنات --- موديول 1: الهوية والوصول والملف الشخصيّ}
\diag{32}{مخطّط المكوّنات --- موديول 2: المنح والطلبات والمستندات}
\diag{33}{مخطّط المكوّنات --- موديول 3: حجز الاستشاريين والتقييمات}
\diag{34}{مخطّط المكوّنات --- موديول 4: المجتمع والدردشة}
\diag{35}{مخطّط المكوّنات --- موديول 5: مركز الموارد والإشعارات}
\diag{36}{مخطّط المكوّنات --- موديول 6: الذكاء الاصطناعيّ وقاعدة المعرفة والمنصّة}
